\documentclass[12pt]{article}
\usepackage[utf8]{inputenc}
\usepackage[T1]{fontenc}
\usepackage[margin=1in]{geometry}
\usepackage{setspace}
\usepackage{xcolor}
\usepackage[colorlinks=true, urlcolor=blue!60!black, linkcolor=blue!60!black]{hyperref}
\setlength{\parindent}{0em}
\setlength{\parskip}{0.6em}
\renewcommand{\baselinestretch}{1.0}
\begin{document}
\begin{center}
{\LARGE \textit{City of Cleburne v. Cleburne Living Center, Inc.}, 473 U.S. 432 (1985)\par}
\vspace{0.5em}
{\large Case Brief --- Two-Page Edition\par}
\end{center}

\section{Memory Jogger}
Intellectual disability is not a quasi-suspect class, but a city fails rational basis review when it denies a group-home permit out of irrational prejudice and unsubstantiated neighborhood fears (\href{https://supreme.justia.com/cases/federal/us/473/432/}{473 U.S. 432}).

\section{Detailed Case Facts}
Cleburne Living Center sought to run a regulated group home for thirteen adults with intellectual disabilities at 201 Featherston Street, in a zone that allowed apartments, boarding houses, dormitories, hospitals, and nursing homes by right. The city deemed it a ``hospital for the feebleminded'' needing a special use permit, and the City Council denied one, citing neighbors' fears, possible harassment by junior high students, a 500-year flood plain, liability, and the home's size.

\section{Procedural History}
The district court upheld the ordinance under rational basis review; the Fifth Circuit reversed, treating intellectual disability as quasi-suspect. The Supreme Court affirmed as to the Featherston home and otherwise vacated.

\section{Judicial Votes}
Unanimous on the judgment invalidating the permit denial; 6--3 on the standard of review. White wrote for six; Stevens, joined by Burger, concurred; Marshall, joined by Brennan and Blackmun, concurred in the judgment in part and dissented in part.

\section{Holding}
Intellectual disability classifications receive rational basis review, but ``mere negative attitudes, or fear'' and ``irrational prejudice'' are not legitimate interests. The permit requirement is invalid as applied.

\section{Analysis of Opinions}
White denied quasi-suspect status because the group's differences are relevant to legitimate goals, remedial legislation shows political responsiveness, and heightened scrutiny would sweep in ``the aging, the disabled, the mentally ill, and the infirm.'' Yet because the zone freely admitted comparable group dwellings, each of the city's reasons failed: neighbor fear cannot be given legal effect, and the flood, liability, and density concerns applied equally to permitted uses. Stevens saw ``a continuum of judgmental responses'' rather than tiers. Marshall called the method ``second order'' rational basis review and, citing a ``lengthy and tragic history'' of eugenics and exclusion, would have applied heightened scrutiny and struck the ordinance on its face.

\section{Examples: Future Applications}
\textbf{Same-Side:} a permit and neighbor-consent rule only for a group home for adults with autism, while student and patient residences are allowed by right, is invalid. \textbf{Opposite-Side:} a cognitive exam for commercial bus drivers that disqualifies some applicants with intellectual disabilities is upheld. \textbf{Fence-Sitter:} a permit rule for all large residential care facilities based on evacuation times, if equally large dormitories are exempt.

\section{Critique}
Scholars call \textit{Cleburne} ``rational basis with bite,'' searching review under a deferential label that leaves lower courts without criteria. Antisubordination critics, echoing Marshall, note that legislative progress cost the group heightened protection, with consequences in \textit{Garrett}; libertarian originalists think the searching inquiry sound but wish it were candid.

\section{Key Quotations}
\begin{quote}
``The Equal Protection Clause of the Fourteenth Amendment commands that no State shall `deny to any person within its jurisdiction the equal protection of the laws,' which is essentially a direction that all persons similarly situated should be treated alike.''
\end{quote}
\noindent --- Justice White, for the Court, 473 U.S. at 439.

\end{document}
